\documentclass[10pt,a4paper]{amsart}
\usepackage[margin=19mm]{geometry}
\usepackage{amsmath,amssymb,mathtools}
\usepackage[scr=rsfs,cal=euler]{mathalfa}
\usepackage{xfrac}
\usepackage[unicode,hidelinks,bookmarks=false]{hyperref}
\newcommand{\ie}{i.e.\ }
\newcommand{\cf}{cf.\ }
\newcommand{\resp}{respectively\ }
\newcommand{\Subsection}{\subsection}
\providecommand{\nobreakdash}{\nobreak\hbox{-}}
\setlength{\emergencystretch}{2em}
\multlinegap0pt
\usepackage{enumerate}
\usepackage{xcolor}
\usepackage[textsize=tiny,disable]{todonotes}
\usepackage{cleveref}

%formatting

\newcommand{\dfn}{:=}
\renewcommand{\emptyset}{\varnothing}
\renewcommand{\setminus}{-}


\newcommand{\red}[1]{{\color{red}#1}}
\newcommand{\blue}[1]{{\color{blue}#1}}
\newcommand{\purple}[1]{{\color{purple}#1}}
\definecolor{mygreen}{rgb}{0, 0.5, 0}
\newcommand{\green}[1]{{\color{mygreen}#1}}
\newcommand{\orange}[1]{{\color{orange}#1}}

\newcommand{\anupam}[1]{\todo{A: #1}}
\newcommand{\gianluca}[1]{\todo{G: #1}}

\newcommand{\toappendix}[1]{\purple{#1}}
% \newcommand{\toappendix}[1]{} %to delete everything inside \toappendix command

%theorems (same as cyclic-bc file)

\newtheorem{thm}{Theorem}
\newtheorem{lem}[thm]{Lemma}
\newtheorem{prop}[thm]{Proposition}
\newtheorem{cor}[thm]{Corollary}
\newtheorem{conj}[thm]{Conjecture}
\newtheorem{fact}[thm]{Fact}
\newtheorem{cond}[thm]{Condition}

\theoremstyle{definition}
\newtheorem{defn}[thm]{Definition}
\newtheorem{exmp}[thm]{Example}
\newtheorem{conv}[thm]{Convention}

\theoremstyle{remark}
\newtheorem{rem}[thm]{Remark}
\newtheorem*{note}{Note}
\newtheorem*{nota}{Notation}


%sets

\newcommand{\Nat}{\mathbb{N}}
\newcommand{\N}{\Nat}
\newcommand{\Bool}{\{0,1 \}}

%relations

\newcommand{\wfpo}{\vartriangleleft}
\newcommand{\wfpoeq}{\trianglelefteq}
\newcommand{\pref}{\prec}
\newcommand{\prefeq}{\preceq}

\newcommand{\permpref}{\subset}
\newcommand{\permprefeq}{\subseteq}

\newcommand{\pp}{\permpref;\permprefeq}

%types
\newcommand{\n}{N}
\newcommand{\bool}{B}
\newcommand{\sq}{\Box}
\newcommand{\sqn}{{\sq\n}}
\newcommand{\sqb}{{\sq\bool}}
\newcommand{\sn}{\sqn}
\newcommand{\lists}[3]{{#2},\overset{#1}{\ldots},{#3} }

%complexity classes

\newcommand{\ptime}{\mathbf{PTIME}}
\newcommand{\pspace}{\mathbf{PSPACE}}
\newcommand{\elementary}{\mathbf{ELEMENTARY}}
\newcommand{\E}{\mathcal{E}}
\newcommand{\Ebin}{\mathcal{E}_{0,1}}

\newcommand{\f}{\mathbf{F}}
\newcommand{\fptime}{\f\ptime}
\newcommand{\fpspace}{\f\pspace}
\newcommand{\felementary}{\f\elementary}
\newcommand{\alogtime}{\mathbf{ALOGTIME}}
\newcommand{\nc}{\mathbf{NC}}
\newcommand{\flinspace}{\mathbf{FLINSPACE}}

%functions

\newcommand{\num}[1]{\underline{#1}}
\renewcommand{\succ}[1]{\mathsf{s}_{#1}}
\newcommand{\pred}{\mathsf{p}}
\newcommand{\floor}[1]{\lfloor {#1} \rfloor}
\newcommand{\hlf}[1]{\floor{#1 /2}}
\newcommand{\cnd}{\mathsf{cond}}
\newcommand{\proj}[2]{\pi^{#1}_{#2}}
\newcommand{\s}[1]{\vert {#1}\vert}
\newcommand{\depth}[1]{\mathsf{depth}{#1}}
\newcommand{\model}[1]{\denot{#1}}
\newcommand{\denot}[1]{f_{#1}}
\newcommand{\ex}{\mathsf{ex}}
\newcommand{\comp}{\mathsf{comp}}
\newcommand{\poly}{\mathrm{poly}}

\newcommand{\sumlen}[1]{|\hspace{-.15em}|#1|\hspace{-.15em}|}

\newcommand{\sexp}{\mathcal{E}}
\newcommand{\cconc}{\mathcal{C}}
\newcommand{\psearch}{\mathcal{B}}
\newcommand{\incrementation}{\mathcal{I}}
\newcommand{\successor}{\mathcal{S}}
\newcommand{\predecessor}{\mathcal{P}}
\newcommand{\completeness}{\mathcal{F}}
\newcommand{\primrec}{\mathcal{R}}
\newcommand{\lengths}{\mathcal{L}}
\newcommand{\unbounded}{\mathcal{U}}
\newcommand{\binarytounary}{\mathcal{N}}
\newcommand{\gfunction}{\mathcal{G}}
\newcommand{\hfunction}{\mathcal{H}}
\newcommand{\saferecursion}{\mathcal{SR}}

\newcommand{\numeral}[1]{\underline{#1}}
\newcommand{\permi}[1]{\vec{#1}}

%schemes

\newcommand{\saferec}{\mathsf{srec}}
\newcommand{\srec}{\saferec}
\newcommand{\saferecwfpo}[1]{\saferec_{#1}}
\newcommand{\srecwfpo}[1]{\saferecwfpo{#1}}

\newcommand{\srecpp}{\srecwfpo{\permpref}}

\newcommand{\safenesrec}{\mathsf{snrec}}
\newcommand{\snrec}{\safenesrec}
\newcommand{\safenesrecwfpo}[1]{\safenesrec_{#1}}
\newcommand{\snrecwfpo}[1]{\safenesrecwfpo{#1}}

\newcommand{\snrecpp}{\snrecwfpo{\permpref}}

\newcommand{\ssrecwfpo}[1]{\mathsf{ssrec}_{#1}}
\newcommand{\ssnrecwfpo}[1]{\mathsf{ssnrec}_{#1}}

\newcommand{\ssrecpp}{\ssrecwfpo{\permpref}}
\newcommand{\ssnrecpp}{\ssnrecwfpo{\permpref}}

%algebras

\newcommand{\bc}{\mathsf{B}}
\newcommand{\bcwfpo}[1]{\bc^{#1}}
\newcommand{\bcnorec}{\bc^{-}}

\newcommand{\nested}{\mathsf{N}}
\newcommand{\nbc}{\nested\bc}
\newcommand{\nbcwfpo}[1]{\nested\bcwfpo{#1}}

\newcommand{\shallow}{\mathsf{S}}
\newcommand{\sbc}{\shallow\bc}
\newcommand{\sbcwfpo}[1]{\sbc^{#1}}

\newcommand{\bcpp}{\bcwfpo{\permpref}}
\newcommand{\nbcpp}{\nbcwfpo{\permpref}}
\newcommand{\sbcpp}{\sbcwfpo\permpref}

%proofs

\newcommand{\seqar}{\Rightarrow}
\newcommand{\der}{\mathcal{D}}
\newcommand{\pder}[1]{\der_{#1}}
\newcommand{\rd}{{R}_\der}
\newcommand{\derrd}{\langle \der, \rd \rangle}
\newcommand{\bud}{\mathit{Bud}}
\newcommand{\nodordereq}{\sqsubseteq}
\newcommand{\nodorder}{\sqsubset}
\newcommand{\close}[1]{{C}_{#1}}
\newcommand{\open}[1]{{O}_{#1}}
 \providecommand\omicron{o}
 \newcommand{\hnu}{\nu_0}

\newcommand{\id}{\mathsf{id}}
\newcommand{\wk}{\mathsf{w}}
\newcommand{\exch}{\mathsf{e}}
\newcommand{\cntr}{\mathsf{c}}
\newcommand{\cut}{\mathsf{cut}}
\newcommand{\pcut}{\mathsf{pcut}}
\newcommand{\boxlef}{\sq_l}
\newcommand{\boxrig}{\sq_r}
\newcommand{\sql}{\boxlef}
\newcommand{\sqr}{\boxrig}
\newcommand{\zero}{0}
\newcommand{\com}{\mathsf{dis}}
\newcommand{\rules}{\mathsf{r}}





%systems

\newcommand{\cyclic}{\mathsf{C}}
\newcommand{\cbc}{\cyclic\bc}
\newcommand{\ncbc}{\cyclic\nested\bc}
\newcommand{\scbc}{\shallow\cbc}
\newcommand{\slr}{\mathsf{SLR}}
\newcommand{\PR}{\mathsf{PR}}
\newcommand{\tgodel}{\mathsf{T}}
\newcommand{\ntgodel}[1]{\tgodel_{#1}}
\newcommand{\ct}{\mathsf{CT}}
\setcounter{section}{4}
\setcounter{thm}{40}
\setcounter{equation}{6}
\let\qoriglem\lem \def\lem{\crefalias{thm}{lem}\qoriglem}
\let\qorigprop\prop \def\prop{\crefalias{thm}{prop}\qorigprop}
\let\qorigcor\cor \def\cor{\crefalias{thm}{cor}\qorigcor}
\let\qorigconj\conj \def\conj{\crefalias{thm}{conj}\qorigconj}
\let\qorigfact\fact \def\fact{\crefalias{thm}{fact}\qorigfact}
\let\qorigcond\cond \def\cond{\crefalias{thm}{cond}\qorigcond}
\let\qorigdefn\defn \def\defn{\crefalias{thm}{defn}\qorigdefn}
\let\qorigexmp\exmp \def\exmp{\crefalias{thm}{exmp}\qorigexmp}
\let\qorigconv\conv \def\conv{\crefalias{thm}{conv}\qorigconv}
\let\qorigrem\rem \def\rem{\crefalias{thm}{rem}\qorigrem}
\title[Soundness of the relativised safe-recursion function algebra]{Cyclic Implicit Complexity: Theorem43}
\author{Gianluca Curzi; Anupam Das}
\date{Source: arXiv:2110.01114v3 (CC BY 4.0)}
\begin{document}
\maketitle
\noindent\emph{Source and licence.} Cyclic Implicit Complexity. Version arXiv:2110.01114v3 (CC BY 4.0). This material is distributed under the Creative Commons Attribution 4.0 International licence, http://creativecommons.org/licenses/by/4.0/, as recorded in the arXiv OAI metadata for this exact version and shown on the abstract page https://arxiv.org/abs/2110.01114v3. Full rights notice: licenses/arxiv-2110.01114-CC-BY-4.0.txt.\par\smallskip
\noindent\textit{Editorial scope.} Counted unit: original Theorem 43 (label thm:fp-soundness), the soundness theorem of the article, with its complete proof. The unit is delivered inside the AMS amsart wrapper; the publisher's acmart class is not used. The article's own macro file is copied into the wrapper, and the section/counter offsets are set so that printed numbering agrees with the published PDF. Counted claim: original Theorem 43 (source label thm:fp-soundness, source0.tex lines 2313-2322), the soundness half of the article's main characterisation, delivered with its complete original proof at source lines 2331-2383 as the single primary proof body. Required context retained uncounted: the whole of Section 5 as printed (source lines 1987-2312 and 2323-2330), that is the section heading, the relativised Bounding Lemma 41 with its complete proof, Remark 42, the definitions of the elementary and binary-elementary functions, the oracle-bound equations (7)-(11) and every cross-reference the counted proof uses. The article sets the section counter to 4 and the shared theorem counter to 40 through editorial \textbackslash{}setcounter instructions placed in the prelude (the source itself never resets them, since the wrapper reproduces only one section), so that all printed numbers coincide with the published PDF: Section 5, Definition/Lemma 41, Remark 42, Theorem 43 and equations (7)-(11). The article's own macro file cbc-macros.tex is copied verbatim except that its \textbackslash{}newcommand\textbackslash{}omicron is written \textbackslash{}providecommand\textbackslash{}omicron because the pinned LaTeX kernel already defines that name; this is an editorial compatibility change to one macro declaration, not to any mathematics. Nothing else is omitted, no source line is dropped, and no artwork is redistributed; the article's own virginialake proof-tree macros are not needed by this unit and are not included.
\section{Characterisations for function algebras}\label{sec:characterization-results-for-function-algebras}
In this section we characterise the complexities  of the function algebras we introduced in the previous section. 
Namely, despite apparently extending $\bc$, $\bcpp$ still contains just the polynomial-time functions, whereas  both  $\nbc$ and $\nbcpp$ are shown to contain just the elementary functions. 
All such results rely on a `bounding lemma' inspired by~\cite{BellantoniCook}.




\subsection{A relativised Bounding Lemma}
\label{sec:bounding-lemma}
Bellantoni and Cook showed in \cite{BellantoniCook} that any function $f(\vec x;\vec y) \in \bc$ satisfies the `poly-max bounding lemma': there is a polynomial $p_f(\vec n)$ such that:\footnote{Recall that, for $\vec x = x_1, \dots, x_n$, we write $|\vec x|$ for $|x_1|, \dots, |x_n|$.}
\begin{equation}
    \label{eq:bc-polymax-eqn}
    |f(\vec x;\vec y)| \leq p_f (|\vec x|) + \max |\vec y| 
\end{equation}
This provided a suitable invariant for eventually proving that all $\bc$-functions were polynomial-time computable.

In this work, inspired by that result, we generalise the bounding lemma to a form suitable to the relativised algebras from the previous section. 
To this end we establish in the next result a sort of `elementary-max' bounding lemma that accounts for the usual poly-max bounding as a special case, by appealing to the notion of (un)nested safe composition.
Both the statement and the proof are quite delicate due to our algebras' formulation using oracles; we must assume an appropriate bound for the oracles themselves, and the various (mutual) dependencies in the statement are subtle.


To state and prove the Bounding Lemma, let us employ the notation $\sumlen{\vec x} \dfn \sum |\vec x|$. 
% \begin{lem}[Bounding Lemma]
% \label{lem:boundinglemma}
% % \marginpar{\footnotesize See proof at \ref{prf:boundinglemma}.}
% Let $f(\vec a)(\vec x;\vec y) \in \nbcpp(\vec a)$, with $\vec a = a_1, \dots, a_k$.
% There is a function $m_f^{\vec c}(\vec x,\vec y)$ with,
% \[
% m_f^{\vec c}(\vec x,\vec y) = 
% e_f(\sumlen{\vec x}) + d_f\sum\vec c + \max |\vec y|
% \]
% for an elementary function $e_f(n)$ and a constant $d_f \geq 1$, such that whenever there are constants $\vec c=c_1, \dots, c_k$ such that,
% \begin{equation}
%     \label{eq:oracle-const-max-bound}
%     |a_i (\vec x_i;\vec y_i)| \leq c_i + d_f \sum_{j\neq i} c_j + \max|\vec y_i|
% \end{equation}
% % for some constants $\vec c =c_1, \ldots, c_k$, then 
% for $1\leq i \leq k$,
% we have:\footnote{To be clear,  we write $|\vec y_i| \leq m_f^{\vec c}(\vec x,\vec y)  $ here as an abbreviation for $\{|y_{ij}| \leq m_f^{\vec c}(\vec x,\vec y)\}_j$.}
% \begin{equation}
%     \label{eq:elem-const-max-bound}
%     |f(\vec a)(\vec x;\vec y)| \leq m_f^{\vec c}(\vec x,\vec y)
% \end{equation}
% \begin{equation}
%     \label{eq:input-bounding-eqn}
%     f(\vec a)(\vec x; \vec y) \ = \ f\left(\lambda \vec x_i.\lambda |\vec y_i| \leq m_f^{\vec c}(\vec x,\vec y) . a_i(\vec x_i;\vec y_i)\right)_i(\vec x;\vec y)
% \end{equation}
% % \anupam{this last equation a little difficult to parse.}
% % where \red{$m_f^{\vec c}(\vec x,\vec y)$ is shorthand notation for:
% % \[
% % e_f(\sumlen{\vec x}) + d_f\sum\vec c + \max |\vec y|
% % \]
% % }
% Moreover, if in fact $f(\vec x;\vec y) \in \bcpp(\vec a) $, then $d_f=1$ and $e_f(n)$ is a polynomial. 
% \end{lem}

\begin{lem}[Bounding Lemma]
\label{lem:boundinglemma}
% \marginpar{\footnotesize See proof at \ref{prf:boundinglemma}.}
Let $f(\vec a)(\vec x;\vec y) \in \nbcpp(\vec a)$, with $\vec a = a_1, \dots, a_k$, and suppose $g_1, \ldots, g_k$ are safe-normal functions. Then there is an elementary function $e_f(n)$ and a constant $d_f \geq 1$ such that whenever there are constants $\vec c=c_1, \dots, c_k$ satisfying, for $1\leq i \leq k$,
\begin{equation}
    \label{eq:oracle-const-max-bound}
    |g_i (\vec u;\vec v)| \leq c_i + \left(d_f \sum_{j\neq i} c_j \right) + \max|\vec v| \qquad \qquad \text{for all inputs } \vec u,\vec v
\end{equation}
% for some constants $\vec c =c_1, \ldots, c_k$, then 
then we have the following, for all inputs $\vec u, \vec v$\footnote{To be clear,  we write $|\vec y| \leq m_f(\vec c,\vec u,\vec v)  $ here as an abbreviation for $\{|y_{j}| \leq m_f(\vec c,\vec u,\vec v)\}_j$.}:
\begin{equation}
    \label{eq:elem-const-max-bound}
    |f(\vec g)(\vec u;\vec v)| \leq m_f(\vec c, \vec u,\vec v)
\end{equation}
\begin{equation}
    \label{eq:input-bounding-eqn}
    f(\vec g)(\vec u; \vec v) \ = \ f\left(\lambda \vec x.\lambda |\vec y| \leq m_f(\vec c, \vec u,\vec v) . g_i(\vec x;\vec y)\right)_{1 \leq i \leq k}(\vec u;\vec v)
\end{equation}
where 
\[
m_f(\vec c,\vec x,\vec y) = 
e_f(\sumlen{\vec x}) + d_f\sum\vec c + \max |\vec y|
\]
% \anupam{this last equation a little difficult to parse.}
% where \red{$m_f^{\vec c}(\vec x,\vec y)$ is shorthand notation for:
% \[
% e_f(\sumlen{\vec x}) + d_f\sum\vec c + \max |\vec y|
% \]
% }
Moreover, if in fact $f(\vec x;\vec y) \in \bcpp(\vec a) $, then $d_f=1$ and $e_f(n)$ is a polynomial. 
\end{lem}

Notice that in the above statement we employ the notation $\lambda \vec x, \lambda|\vec y|\leq n$, where $\vec y=y_1, \ldots, y_k$, for the guarded abstraction defined by:
\[
(\lambda \vec x, \lambda|\vec y|\leq n. h(\vec x; \vec y))(\vec w; \vec z):= \begin{cases}
    h(\vec w; \vec z) &\text{if } |z_j|\leq n \text{ for all }1 \leq j \leq k\\
    0 &\text{otherwise.}
\end{cases}
\]


Unwinding the statement above, note that $e_f$ and $d_f$ depend \emph{only} on the function $f$ itself, not on the constants $\vec c$ given for the (mutual) oracle bounds in \Cref{eq:oracle-const-max-bound}.
This is crucial for the proof, namely in the case when $f$ is defined by recursion, substituting different values for $\vec c$ during an inductive argument.
% This independency of $e_f$ and $d_f$ from $\vec c$ will be made explicit in their definitions in the proof that follows.

While the role of the elementary bounding function $e_f$ is a natural counterpart of $p_f$ in Bellantoni and Cook's bounding lemma, cf.~\Cref{eq:bc-polymax-eqn}, the role of $d_f$ is perhaps slightly less clear.
Intuitively, $d_f$ represents the amount of `nesting' in the definition of $f$, increasing whenever oracle calls are substituted into arguments for other oracles. 
Hence, if $f$ uses only unnested Safe Composition, then $d_f=1$ as required. 
In fact, it is only important to distinguish whether $d_f=1$ or not, since $d_f$ forms the base of an exponent for defining $e_f$ when $f$ is defined by safe recursion.

Finally let us note that~\Cref{eq:elem-const-max-bound,eq:input-bounding-eqn} are somewhat dual:  while the former bounds the \emph{output} of a function (serving as a modulus of \emph{growth}), the latter bounds the \emph{inputs} (serving as a modulus of \emph{continuity}).


\begin{rem}
Let us also point out that we may relativise the statement of the lemma to any set of oracles including those which $f(\vec x;\vec y)$ is over.
In particular, if $f(\vec x;\vec y)$ is over no oracles, then we may realise \Cref{eq:oracle-const-max-bound} vacuously by choosing $\vec a = \emptyset$ and we would obtain that $|f(\vec x;\vec y)| \leq e_f(\sumlen{\vec x}) + \max |\vec y|$.
More interestingly, in the case when $f(\vec x;\vec y)$ is just, say, $a_i(\vec x;\vec y)$, we may choose to set $\vec a = a_i$ or $\vec a = a_1,\dots, a_n$ in the above lemma, yielding different bounds in each case.
We shall exploit this in inductive hypotheses  in the proof that follows (typically when we write `WLoG').
\end{rem}

\begin{proof}[Proof of~\Cref{lem:boundinglemma}.]
\label{prf:boundinglemma}
We prove~\Cref{eq:elem-const-max-bound} and~\Cref{eq:input-bounding-eqn} by induction on the definition of $f(\vec x;\vec y)$, always assuming that we have $\vec c$ satisfying \Cref{eq:oracle-const-max-bound}.

Throughout the argument we shall actually construct $e_f( n)$ that are \emph{monotone} elementary functions (without loss of generality) and exploit this invariant.
In fact $e_f( n)$ will always be generated by composition from $0, \succ{}, +, \times,  x^y$ and projections.
If $f(\vec x;\vec y) \in \bcpp(\vec a)$ then $e_f(n)$ will be in the same algebra without exponentiation, $ x^y$, i.e.\ it will be a polynomial with only non-negative coefficients.
This property will be made clear by the given explicit definitions of $e_f( n)$ throughout the argument.


Let us start with~\Cref{eq:elem-const-max-bound}. If $f(\vec x;\vec y)$ is an initial function then it suffices to set $e_f(n) \dfn 1 +  n$ and $d_f\dfn 1$.

If $f(\vec x;\vec y) = a_i(\vec x;\vec y)$ then it suffices to set $e_f(n) \dfn 0 $ and $d_f \dfn 1$.

If $f(\vec x;\vec y) = h(\vec x;\vec y, g(\vec x;\vec y))$, let $e_h, e_g, d_h, d_g$ be obtained from the inductive hypothesis.
We have,
\[
\arraycolsep=2pt
\begin{array}{rcl}
    |f(\vec x;\vec y)| & = & |h(\vec x; \vec y, g(\vec x; \vec y))| \\
    & \leq & e_h(\sumlen{\vec x}) + d_h\sum\vec c + \max(|\vec y|, |g(\vec x;\vec y)|) \\
    & \leq & e_h(\sumlen{\vec x}) + d_h\sum\vec c + e_g(\sumlen{\vec x}) + d_g\sum \vec c + \max|\vec y|
    % \\
    % & \leq & (e_h(\sumlen{\vec x}) + e_g(\sumlen{\vec x})) + (d_h+d_g)\sum \vec c + \max |\vec y
\end{array}
\]
so we may set $e_f(n) \dfn e_h(n) + e_g(n)$ and $d_f\dfn d_h+d_g$.

For the `moreover' clause, note that if $f(\vec x;\vec y)$ uses only unnested Safe Composition, then one of $g$ or $h$ does not have oracles, and so we can assume WLoG that either the term $d_h\sum \vec c$ or the term $d_g\sum \vec c$ above does not occur, and we set $d_f \dfn d_g$ or $d_f\dfn d_h$, respectively. In either case we obtain $d_f=1$ by the inductive hypothesis, as required.

If $f(\vec x;\vec y) = h(\vec x, g(\vec x;);\vec y)$, let $e_h,e_g,d_h, d_g$ be obtained from the inductive hypothesis. Note that, by definition of Safe Composition along a normal parameter, we must have that $g$ has no oracles, and so in fact $|g(\vec x;)| \leq e_g(\sumlen{\vec x})$.
We thus have,
\[
\arraycolsep=2pt
\begin{array}{rcl}
    |f(\vec x;\vec y)| & = & |h(\vec x, g(\vec x;);\vec y)| \\
     & \leq & e_h(\sumlen{\vec x}, |g(\vec x;)|) + d_h\sum \vec c + \max|\vec y| \\
     & \leq & e_h(\sumlen{\vec x}, e_g(\sumlen{\vec x})) + d_h\sum \vec c + \max|\vec y| 
\end{array}
\]
so we may set $e_f(n)\dfn e_h(n, e_g(n))$ and $d_f \dfn d_h$.
% 
For the `moreover' clause, note that by the inductive hypothesis $e_h,e_g$ are polynomials and $d_h=1$, so indeed $e_f$ is a polynomial and $d_f=1$.

Finally, if $f(\vec x;\vec y) = h(\lambda \vec u \permpref \vec x, \lambda \vec v. f(\vec u;\vec v) ) (\vec x;\vec y)$, let $e_h,d_h$ be obtained from the inductive hypothesis.
We claim that it suffices to set $d_f \dfn d_h$ and $e_f(n) \dfn nd_h^n e_h(n)$.
Note that, for the `moreover' clause, if $d_h=1$ then also $d_h^n = 1$ and so indeed $e_f(n)$ is a polynomial if $d_h=1$ and $e_h(n)$ is a polynomial.

First, let us calculate the following invariant, for $n>0$:
\begin{equation*}
    \arraycolsep=2pt
    \begin{array}{rcll}
    e_f(n) & = & nd_h^n e_h(n) &  \\
        & = & d_h^n e_h(n) + (n-1)d_h^n e_h(n) & \\
        & = & d_h^n e_h(n) + d_h (n-1) d_h^{n-1}e_h(n) \\
        & \geq & e_h(n) + d_h (n-1) d_h^{n-1}e_h(n) & \text{$d_h\geq 1$} \\
        & \geq & e_h(n) + d_h (n-1) d_h^{n-1}e_h(n-1) \  &  \text{$e_h$~monotone} \\
        & \geq & e_h(n) + d_h e_f(n-1) & \text{def.~of $e_f$}
\end{array}
\end{equation*}
Hence:
\begin{equation}
 \label{eq:bnd-lem-rec-fn-inv}
     e_f(n) \geq  e_h(n) + d_h e_f(n-1) 
\end{equation}




Now, to show that \Cref{eq:elem-const-max-bound} bounds the $f,e_f,d_f$ at hand, we proceed by a sub-induction on $\sumlen{\vec x}$.
For the base case, when $\sumlen{\vec x} = 0 $ (and so, indeed, $\vec x = \vec 0$), note simply that $\lambda \vec u\permpref \vec x, \lambda \vec v . f(\vec u;\vec v)$ is the constant function $0$, and so we may appeal to the main inductive hypothesis for $h(a)$ setting the corresponding constant $c$ for $a$ to be $0$ to obtain,
\[
\arraycolsep=2pt
\begin{array}{rcl}
    |f(\vec 0;\vec y)| & = & |h(0)(\vec 0;\vec y)| \\
        & \leq & e_h(0) + d_h\sum\vec c + \max |\vec y| \\
        & \leq & e_f(0) + d_f\sum\vec c + \max |\vec y|
\end{array}
\]
as required.
For the sub-inductive step,  let $\sumlen{\vec x} > 0$. 
Note that, whenever $\vec u \permpref \vec x$ we have $\sumlen{\vec u} < \sumlen{\vec x} $ and so, by the sub-inductive hypothesis and monotonicity of $e_f$ we have:
\[
|f(\vec u;\vec v)| \leq e_f(\sumlen {\vec x} -1) + d_f \sum \vec c + \max |\vec v|
\]
Now we may again appeal to the main inductive hypothesis for $h(a)$ by setting $c= e_f(\sumlen{\vec x}  -1)$ to be the corresponding constant for $a = \lambda \vec u \permpref \vec x, \lambda \vec v. f(\vec u;\vec v)$. 
We thus obtain:
\[
\arraycolsep=2pt
\begin{array}{rclll}
     |f(\vec x;\vec y)| & = & |h(\lambda \vec u \permpref \vec x, \lambda \vec v . f(\vec u;\vec v)) (\vec x;\vec y)| \\
         & \leq & e_h (\sumlen{\vec x}) + d_h c + d_h\sum\vec c + \max |\vec y| & & \text{main IH} \\
         & \leq & ( e_h (\sumlen{\vec x}) + d_h e_f(\sumlen{\vec x} -1) ) +\\
         & &  + d_h\sum\vec c + \max |\vec y| & & \text{def.~of $c$}\\
         & \leq & e_f(\sumlen{\vec x}) + d_h \sum \vec c + \max|\vec y| && \eqref{eq:bnd-lem-rec-fn-inv} \\ & \leq & e_f(\sumlen{\vec x}) + d_f \sum \vec c + \max|\vec y| && \text{$d_f = d_h$}
\end{array}
\]
Let us now prove~\Cref{eq:input-bounding-eqn}, and let $e_f$ and $d_f$ be constructed by induction on $f$ as above. We proceed again by induction on the definition of $f(\vec a)(\vec x;\vec y)$, always making explicit the oracles of a function.
% Let us temporarily write $m_f(\vec x,\vec y) \dfn e_f(\sumlen{\vec x}) + d_f\sum\vec c + \max |\vec y|$.
\gianluca{Is it better to define $e_f$ and $d_f$ in advance before proving the two inequations?}

The initial functions and oracle calls are immediate, due to the `$\max |\vec y|$' term in \Cref{eq:input-bounding-eqn}. 

If $f(\vec a)(\vec x; \vec y) = h(\vec a)(\vec x; \vec y, g(\vec a)(\vec x; \vec y))$ then, by the inductive hypothesis for $h(\vec a)$, any oracle call from $h(\vec a)$ only takes safe inputs of lengths:
    \[
    \begin{array}{rll}
        \leq & e_h(\sumlen{\vec x}) + d_h \sum \vec c + \max (|\vec y|, |g(\vec a)(\vec x;\vec y)|) \\
        \leq & e_h(\sumlen{\vec x})+ d_h \sum \vec c + e_g(\sumlen{\vec x}) + d_g\sum \vec c + \max |\vec y| & \text{\eqref{eq:oracle-const-max-bound}} \\
        \leq &  (e_h(\sumlen{\vec x} + e_g(\sumlen{\vec x})) + (d_h +d_g)\sum\vec c + \max |\vec y|\\
         \leq & e_f(\sumlen{\vec x}) + d_f\sum\vec c + \max|\vec y|
    \end{array}
    \]
    Note that any oracle call from $g(\vec a)$ will still only take safe inputs of lengths $\leq e_g(\sumlen{\vec x}) + d_g\sum \vec c + \max |\vec y|$, by the inductive hypothesis, and $e_g$ and $d_g$ are bounded above by $e_f$ and $d_f $ respectively.
    
If $f(\vec a)(\vec x;\vec y) = h(\vec a)(\vec x, g(\emptyset)(\vec x;);\vec y)$ then, by the inductive hypothesis, any oracle call will only take safe inputs of lengths:
\[
\arraycolsep=2pt
\begin{array}{rlll}
    \leq & e_h(\sumlen{\vec x} + |g(\emptyset)(\vec x;)|) + d_h\sum\vec c + \max |\vec y| \\
    \leq & e_h(\sumlen{\vec x} + e_g(\sumlen{\vec x})) + d_h\sum\vec c + \max |\vec y| & & \text{\Cref{lem:boundinglemma}} \\
    \leq & e_f(\sumlen{\vec x}) + d_f\sum \vec c + \max |\vec y|
\end{array}
\]

Last, suppose $f(\vec a)(\vec x;\vec y)= h(\vec a, \lambda \vec u\permpref \vec x, \lambda \vec v . f(\vec a)(\vec u;\vec v))(\vec x;\vec y)$. 
We proceed by a sub-induction on $\sumlen{\vec x}$.
Note that, since $\vec u \permpref \vec x \implies \sumlen{\vec u} <\sumlen{\vec x}$, we immediately inherit from the inductive hypothesis the appropriate bound on safe inputs for oracle calls from $\lambda \vec u\permpref \vec x, \lambda \vec v. f(\vec a)(\vec u;\vec v)$.

Now, recall from the Bounding \Cref{lem:boundinglemma}, whenever $\vec u \permpref \vec x$ (and so $\sumlen{\vec u} < \sumlen{\vec x}$), we have $|f(\vec u;\vec v)| \leq e_f(\sumlen{\vec x}-1) + d_f\sum \vec c + \max |\vec v|$.
So by setting $c = e_f(\sumlen{\vec x}-1)$ in the inductive hypothesis for $h(\vec a,a)$, with $a = \lambda \vec u\permpref \vec x, \lambda \vec v. f(\vec a)(\vec u;\vec v)$, any oracle call from $h(\vec a,a)$ will only take safe inputs of lengths:
\[
\arraycolsep=2pt
\begin{array}{rlll}
    \leq & e_h(\sumlen{\vec x}) + d_h e_f(\sumlen{\vec x}-1) + d_h\sum \vec c + \max |\vec y| \\
    \leq & e_f(\sumlen{\vec x}) + d_f \sum \vec c + \max |\vec y| & & \text{\eqref{eq:bnd-lem-rec-fn-inv}}
\end{array}
\]
% Now, for an oracle call from $\lambda \vec u\permpref \vec x, \lambda \vec v. f(\vec a)(\vec u;\vec v)$ we may simply apply the same reasoning again, essentially proceeding by induction on $\sumlen{\vec x}$.
This completes the proof.  
\end{proof}



% Before moving on to our soundness results, let us first present an important consequence of the bounding lemma that will be important for characterising $\nbcpp$.

% \begin{prop}
% % \label{input-elem-bound}
% Let $f(\vec a)(\vec x; \vec y) \in \nbcpp (\vec a) $ and let $e_f$ and $d_f$ be as constructed in the proof of \Cref{lem:boundinglemma} and write:
% \[
% m_f^{\vec c}(\vec x,\vec y) \ \dfn \ e_f(\sumlen{\vec x}) + d_f\sum\vec c + \max |\vec y|
% \]
% If there are constants $\vec c$ such that $\vec a$ satisfy \Cref{eq:oracle-const-max-bound},
% then:\footnote{To be clear, here we write $|\vec y_i| \leq m_f^{\vec c}(\vec x,\vec y)  $ here as an abbreviation for $\{|y_{ij} \leq m_f^{\vec c}(\vec x,\vec y)\}_j$.}
% \begin{equation}
%     \label{eq:input-bounding-eqn}
%     f(\vec a)(\vec x; \vec y) \ = \ f(\lambda \vec x_i.\lambda |\vec y_i| \leq m_f^{\vec c}(\vec x,\vec y) . a_i(\vec x_i;\vec y_i))_i(\vec x;\vec y)
% \end{equation}
% \end{prop}





% \toappendix{\begin{proof}
% We proceed by induction on the definition of $f(\vec a)(\vec x;\vec y)$, always making explicit the oracles of a function.
% % Let us temporarily write $m_f(\vec x,\vec y) \dfn e_f(\sumlen{\vec x}) + d_f\sum\vec c + \max |\vec y|$.

% The initial functions and oracle calls are immediate, due to the `$\max |\vec y|$' term in \Cref{eq:input-bounding-eqn}. 

% If $f(\vec a)(\vec x; \vec y) = h(\vec a)(\vec x; \vec y, g(\vec a)(\vec x; \vec y))$ then, by the inductive hypothesis for $h(\vec a)$, any oracle call from $h(\vec a)$ only take safe inputs of lengths:
%     \[
%     \begin{array}{rll}
%         \leq & e_h(\sumlen{\vec x}) + d_h \sum \vec c + \max (|\vec y|, |g(\vec a)(\vec x;\vec y)|) \\
%         \leq & e_h(\sumlen{\vec x})+ d_h \sum \vec c + e_g(\sumlen{\vec x}) + d_g\sum \vec c + \max |\vec y| & \text{by Bounding \Cref{lem:boundinglemma}} \\
%         \leq &  (e_h(\sumlen{\vec x} + e_g(\sumlen{\vec x})) + (d_h +d_g)\sum\vec c + \max |\vec y|\\
%          \leq & e_f(\sumlen{\vec x}) + d_f\sum\vec c + \max|\vec y|
%     \end{array}
%     \]
%     Note that any oracle call from $g(\vec a)$ will still only take safe inputs of lengths $\leq e_g(\sumlen{\vec x}) + d_g\sum \vec c + \max |\vec y|$, by the inductive hypothesis, and $e_g$ and $d_g$ are bounded above by $e_f$ and $d_f $ respectively.
    
% If $f(\vec a)(\vec x;\vec y) = h(\vec a)(\vec x, g(\emptyset)(\vec x;);\vec y)$ then, by the inductive hypothesis, any oracle call will only take safe inputs of lengths:
% \[
% \begin{array}{rll}
%     \leq & e_h(\sumlen{\vec x} + |g(\emptyset)(\vec x;)|) + d_h\sum\vec c + \max |\vec y| \\
%     \leq & e_h(\sumlen{\vec x} + e_g(\sumlen{\vec x})) + d_h\sum\vec c + \max |\vec y| & \text{by Bounding \Cref{lem:boundinglemma}} \\
%     \leq & e_f(\sumlen{\vec x}) + d_f\sum \vec c + \max |\vec y|
% \end{array}
% \]

% Finally suppose $f(\vec a)(\vec x;\vec y)= h(\vec a, \lambda \vec u\permpref \vec x, \lambda \vec v . f(\vec a)(\vec u;\vec v))(\vec x;\vec y)$. 
% We proceed by a sub-induction on $\sumlen{\vec x}$.
% Note that, since $\vec u \permpref \vec x \implies \sumlen{\vec u} <\sumlen{\vec x}$, we immediately inherit from the inductive hypothesis the appropriate bound on safe inputs for oracle calls from $\lambda \vec u\permpref \vec x, \lambda \vec v. f(\vec a)(\vec u;\vec v)$.

% Now, recall from the Bounding \Cref{lem:boundinglemma}, whenever $\vec u \permpref \vec x$ (and so $\sumlen{\vec u} < \sumlen{\vec x}$), we have $|f(\vec u;\vec v)| \leq e_f(\sumlen{\vec x}-1) + d_f\sum \vec c + \max |\vec v|$.
% So by setting $c = e_f(\sumlen{\vec x}-1)$ in the inductive hypothesis for $h(\vec a,a)$, with $a = \lambda \vec u\permpref \vec x, \lambda \vec v. f(\vec a)(\vec u;\vec v)$, any oracle call from $h(\vec a,a)$ will only take safe inputs of lengths:
% \[
% \begin{array}{rll}
%     \leq & e_h(\sumlen{\vec x}) + d_h e_f(\sumlen{\vec x}-1) + d_h\sum \vec c + \max |\vec y| \\
%     \leq & e_f(\sumlen{\vec x}) + d_f \sum \vec c + \max |\vec y| & \text{by \Cref{eq:bnd-lem-rec-fn-inv}}
% \end{array}
% \]
% % Now, for an oracle call from $\lambda \vec u\permpref \vec x, \lambda \vec v. f(\vec a)(\vec u;\vec v)$ we may simply apply the same reasoning again, essentially proceeding by induction on $\sumlen{\vec x}$.
% This completes the proof.
% \end{proof}
% }

\subsection{Soundness results}
In this subsection we  show that the function algebras $\bcpp$ and $\nbcpp$ (as well as $\nbc$)  capture precisely the classes  $\fptime$ and $\felementary$, respectively. We start with $\bcpp$:


\begin{thm}
\label{thm:fp-soundness}
Suppose $\vec a$ satisfies \Cref{eq:oracle-const-max-bound} for some constants $\vec c$.
We have the following:
\begin{enumerate}
    \item\label{item:soundness-fp} If $f(\vec x;\vec y) \in \bcpp (\vec a)$ then $f(\vec x,\vec y) \in \fptime(\vec a)$.
    % \item\label{item:soundness-fpspace} if $f(\vec x;\vec y) \in \sbcpp(\vec a)$ then $f (\vec x,\vec y) \in \fpspace(\vec a)$, as long as each $d_i = 1$.
    \item\label{item:soundness-felementary} If $f(\vec x;\vec y) \in \nbcpp (\vec a)$ then $f(\vec x,\vec y) \in \felementary(\vec a)$.
\end{enumerate}
\end{thm}



Note in particular that, for $f(\vec x;\vec y)$ in $\bcpp$ or $\nbcpp$, i.e.\ not using any oracles, we immediately obtain membership in $\fptime$ or $\felementary$, respectively.
However, the reliance on intermediate oracles during a function definition causes some difficulties that we must take into account.
At a high level, the idea is to use the Bounding Lemma   (namely \Cref{eq:input-bounding-eqn}) to replace certain oracle calls with explicit appropriately bounded functions computing their graphs. 
From here we compute $f(\vec x;\vec y)$ by a sort of `course-of-values' recursion on $\permpref$, storing previous values in a lookup table.
In the case of $\bcpp$, it is important that this table has polynomial-size, since there are only $m! \prod |\vec x|$ permutations of prefixes of a list $\vec x= x_1, \dots, x_m$ (which is a polynomial of degree $m$).


\begin{proof}[Proof of~\Cref{thm:fp-soundness}.]
We proceed by induction on the definition of $f(\vec x;\vec y)$.

Each initial function is polynomial-time computable, and each (relativised) complexity class considered is under composition, so it suffices to only consider the respective recursion schemes.
We shall focus first on the case of $\bcpp (\vec a)$, \Cref{item:soundness-fp} above, so that $e_f$ is a polynomial and $d_f=1$.



Suppose we have $h(a)(\vec x;\vec y) \in \bcpp(a, \vec a)$ and let:
\[
f(\vec x;\vec y) = h(\lambda \vec u \subset \vec x, \lambda \vec v \subseteq \vec y. f(\vec u; \vec v))(\vec x;\vec y)
\]
We start by making some observations:
\begin{enumerate}
    \item\label{item:poly-growth-rate} First, note that $|f(\vec x;\vec y)| \leq e_f(|\vec x|) +d_f\sum \vec c + \max |\vec y| $, by the Bounding \Cref{lem:boundinglemma}, and so $|f(\vec x;\vec y)|$ is polynomial in $|\vec x,\vec y|$.
%     \begin{itemize}
% \item $|f(\vec x;\vec y)| \leq e_f(|\vec x|) +d_f\sum \vec c + \max |\vec y| $ by the Bounding \Cref{lem:boundinglemma}.
%     \item $h(a)(\vec x;\vec y) \in \fptime(a,\vec a)$, by the inductive hypothesis, whenever $|a(\vec x;\vec y)| \leq c + d_f \sum \vec c + \max |\vec y|$, cf.~\Cref{eq:oracle-const-max-bound}.
%     In particular, by the above point, this holds when we set $a = \lambda \vec u \permpref \vec x, \lambda \vec v \permprefeq \vec y . f(\vec u; \vec v)$, by setting $c = e_f(\sumlen{\vec x}-1)$.
% \end{itemize}
\item\label{item:poly-many-pp} Second, note that the set $[\vec x;\vec y] \dfn \{(\vec u,\vec v) \ | \ \vec u \permpref \vec x , \vec v \permprefeq \vec y\} $ has size polynomial in $|\vec x,\vec y|$:
\begin{itemize}
    \item write $\vec x = x_1 , \dots, x_m$ and $\vec y = y_1, \dots, y_n$.
    \item Each $x_i$ and $y_j$ have only linearly many prefixes, and so there are at most $|x_1|\cdot \cdots \cdot |x_m||y_1| \cdot \cdots \cdot |y_n| \leq \sumlen{\vec x,\vec y}^{m+n}$ many choices of prefixes for all the arguments $\vec x, \vec y$. 
    (This is a polynomial since $m$ and $n$ are global constants).
    \item Additionally, there are $m!$ permutations of the arguments $\vec x$ and $n!$ permutations of the arguments $\vec y$.
    Again, since $m$ and $n$ are global constants, we indeed have
    $|[\vec x;\vec y]| = O(\sumlen{\vec x, \vec y}^{m+n})$, which is polynomial in $|\vec x,\vec y|$.
\end{itemize}
\end{enumerate}

We describe a polynomial-time algorithm for computing $f(\vec x;\vec y)$ (over oracles $\vec a$) by a sort of `course-of-values' recursion on the order $\permpref \times \permprefeq $ on $[\vec x;\vec y]$.

First, for convenience, temporarily extend $\permpref \times \permprefeq$ to a total well-order on $[\vec x;\vec y]$, and write $S$ for the associated successor function.
    Note that $S$ can be computed in polynomial-time from $[\vec x;\vec y]$.


Define $F(\vec x,\vec y) \dfn \langle f(\vec u;\vec v)\rangle_{\vec u \permpref \vec x, \vec v \permprefeq \vec y}$, i.e.\ it is the graph of $\lambda \vec u \permpref \vec x, \lambda \vec v \permprefeq \vec y . f(\vec u;\vec v)$ that we shall use as a `lookup table'.
    Note that $|F(\vec x,\vec y)| $ is polynomial in $|\vec x,\vec y|$ by \Cref{item:poly-growth-rate} and \Cref{item:poly-many-pp} above.
    Now, we can write:\footnote{Here, as abuse of notation, we are now simply identifying $F(\vec x;\vec y)$ with $\lambda \vec u \permpref \vec x, \lambda \vec v \permprefeq \vec y . f(\vec x;\vec y)$.}
    \[
    \arraycolsep=2pt
    \begin{array}{rcl}
        F(S(\vec x,\vec y)) & =&  \langle f(S(\vec x,\vec y)), F(\vec x,\vec y) \rangle  \\
        & = & \langle h(F(\vec x ,\vec y))(\vec x;\vec y),F(\vec x,\vec y)\rangle
    \end{array}
    \]
    Again by \Cref{item:poly-many-pp} (and since $F$ is polynomially bounded), this recursion terminates in polynomial-time.
    We may now simply calculate $f(\vec x;\vec y)$ as $h(F(\vec x,\vec y))(\vec x;\vec y)$.
    
The argument for $\nbcpp$ is similar, though we need not be as careful about computing the size of the lookup tables ($F$ above) for recursive calls.
The key idea is to use the Bounding Lemma (\Cref{eq:elem-const-max-bound}) to bound the safe inputs of recursive calls so that we can adequately store the lookup table for previous values.
\end{proof}

\begin{thebibliography}{99}
\bibitem[6]{BellantoniCook} Stephen Bellantoni and Stephen Cook, ``A New Recursion-Theoretic Characterization of the Polytime Functions (Extended Abstract)'', in \emph{Proceedings of the Twenty-Fourth Annual ACM Symposium on Theory of Computing} (Victoria, British Columbia, Canada), STOC '92, Association for Computing Machinery, New York, NY, USA, 283--293, DOI 10.1145/129712.129740.
\end{thebibliography}
\end{document}
